\section{总结与延伸}

\subsection{讲者的核心总结}

Chris Manning 在本讲中建立了从词向量到神经网络分类器的完整认知链：

\begin{enumerate}
    \item \textbf{优化基础}：梯度下降（特别是 SGD）是训练所有神经网络模型的基础
    \item \textbf{Word2Vec 深入}：Skip-gram 与 CBOW 两种架构，负采样作为高效训练技巧
    \item \textbf{计数方法与 GloVe}：共现矩阵 + SVD 降维是另一种学习词向量的路径，GloVe 统一了两种方法
    \item \textbf{词向量评估}：内在评估（类比、相似度）与外在评估（NER 等下游任务）
    \item \textbf{词义问题}：多义词向量是各义项的加权叠加，可通过稀疏编码恢复
    \item \textbf{神经网络入门}：从 logistic 回归到多层网络，非线性激活函数是深度的关键
\end{enumerate}

\subsection{全课知识图谱}

\begin{center}
\begin{tikzpicture}[node distance=0.8cm, every node/.style={draw, rounded corners, minimum width=3cm, minimum height=0.7cm, font=\small}]
\node (opt) {优化：GD / SGD};
\node (w2v) [below=of opt] {Word2Vec：Skip-gram + 负采样};
\node (count) [below=of w2v] {计数方法：共现矩阵 / SVD / GloVe};
\node (eval) [below=of count] {评估：内在 vs 外在};
\node (sense) [below=of eval] {词义：多义性与叠加};
\node (nn) [below=of sense] {神经网络分类器};
\draw[->, thick] (opt) -- (w2v);
\draw[->, thick] (w2v) -- (count);
\draw[->, thick] (count) -- (eval);
\draw[->, thick] (eval) -- (sense);
\draw[->, thick] (sense) -- (nn);
\node[draw=none, right=1cm of opt, text width=5cm, font=\scriptsize] {损失函数、梯度、学习率、mini-batch};
\node[draw=none, right=1cm of w2v, text width=5cm, font=\scriptsize] {中心词/上下文词向量、sigmoid、稀疏梯度};
\node[draw=none, right=1cm of count, text width=5cm, font=\scriptsize] {共现概率比率、对数双线性模型};
\node[draw=none, right=1cm of eval, text width=5cm, font=\scriptsize] {类比、相似度、NER};
\node[draw=none, right=1cm of sense, text width=5cm, font=\scriptsize] {Superposition、稀疏编码、上下文向量};
\node[draw=none, right=1cm of nn, text width=5cm, font=\scriptsize] {隐藏层、非线性、交叉熵损失};
\end{tikzpicture}
\end{center}

\subsection{关键 Takeaways}

\begin{importantbox}{五条核心原则}
\begin{enumerate}
    \item \textbf{分布式假设是基石}：词的含义由其上下文决定——无论是 Word2Vec 的预测方法还是 GloVe 的计数方法，都是这一假设的不同实现
    \item \textbf{向量空间中的方向编码语义}：词向量不仅捕捉相似性，还通过向量差编码语义关系（类比性质）
    \item \textbf{预测与计数殊途同归}：Word2Vec 和 GloVe 看似不同，实际上都在对共现统计进行建模
    \item \textbf{一词多义不是障碍}：静态词向量通过叠加编码多个义项，上下文词向量则彻底解决了这一问题
    \item \textbf{非线性是深度学习的灵魂}：没有非线性激活函数，再多的层也只是线性变换
\end{enumerate}
\end{importantbox}

\subsection{拓展阅读}

\begin{itemize}
    \item Mikolov et al., 2013: ``Efficient Estimation of Word Representations in Vector Space'' — Word2Vec 原始论文
    \item Mikolov et al., 2013: ``Distributed Representations of Words and Phrases and their Compositionality'' — 负采样等技巧
    \item Pennington, Socher, Manning, 2014: ``GloVe: Global Vectors for Word Representation'' — \url{https://nlp.stanford.edu/projects/glove/}
    \item Levy \& Goldberg, 2014: ``Neural Word Embedding as Implicit Matrix Factorization'' — 证明 Word2Vec 与矩阵分解的等价性
    \item Arora et al., 2018: ``Linear Algebraic Structure of Word Senses, with Applications to Polysemy'' — 从单一词向量恢复义项
    \item CS224N 课程主页：\url{https://web.stanford.edu/class/cs224n/}
\end{itemize}

\end{document}
